\chapter{Aluminium}

\section{Introduction}

Aluminium is the world's most widely used non-ferrous base metal and, despite the energy intensity
of its primary production, it plays a central enabling role in the twin green and digital transitions.
This synthesis is based on a presentation delivered in March 2026 by Christian Leroy on behalf of
European Aluminium, the trade association founded in 1981 that represents the entire European
aluminium value chain.

\section{European Aluminium: Key Facts and the Value Chain}

European Aluminium represents more than 100 member companies operating approximately 600
plants across 30 European countries (EU, EFTA, UK, Turkey, Bosnia-Herzegovina). The industry
generates an annual turnover of \texteuro40 billion and supports over one million direct and indirect
high-value jobs across the continent.

The aluminium value chain spans several distinct stages:

\begin{enumerate}[label=\arabic*.]
    \item \textbf{Bauxite mining} — extraction of the primary ore, predominantly outside Europe.
    \item \textbf{Alumina production} — chemical refining of bauxite into aluminium oxide
          ($\text{Al}_2\text{O}_3$) via the Bayer process.
    \item \textbf{Primary smelting} — electrolytic reduction of alumina to metallic aluminium.
    \item \textbf{Semi-fabrication} — rolling, extrusion, and casting of aluminium into
          intermediate products.
    \item \textbf{Product manufacturing} — transformation of semi-products into finished goods.
    \item \textbf{Use phase} — deployment of aluminium products across end-use sectors.
    \item \textbf{Recycling} — collection, shredding, sorting, and remelting of end-of-life
          aluminium scrap.
\end{enumerate}

Europe currently produces approximately \textbf{6\%} of global primary aluminium, yet accounts for
roughly \textbf{15\%} of global secondary (recycled) aluminium production — a figure that
underscores the continent's comparative strength in the recycling segment. Notably, 51\% of
European aluminium production already comes from recycled sources, and 90\% of aluminium used
in construction and automotive applications is recycled at end of life.

\section{Global and European Market Figures}

Global primary aluminium production has grown dramatically since the 1950s, rising from a few
thousand kilotonnes per year to well above 60{,}000 kilotonnes by the late 2010s. This growth has
been driven by a succession of demand shifts: use in the public sector (1950s), beverage cans
(1970s), automotive applications (1990s), and most recently the green energy transition.

China is by far the dominant producer, accounting for approximately \textbf{59\%} of global
primary output as of the most recent data, compared to 27\% for the rest of the world, 9\% for the
Gulf Cooperation Council states, 4\% for Europe, and roughly 1\% for the United States. This
concentration of Chinese capacity is described in the presentation as a \textit{looming threat} to
global markets due to potential oversupply and price distortion.

From the European perspective, the sourcing picture has shifted markedly since 1980. By the early
2020s, \textbf{54\%} of aluminium consumed in European semi-fabrication was met through net
imports, \textbf{39\%} through total recycling, and only \textbf{7\%} through domestic primary
production — reflecting both the decline of European smelting capacity and the growth of
secondary production.

\section{Main Aluminium Uses}

European aluminium end-use is dominated by five sectors:

\begin{center}
\begin{tabular}{lc}
\toprule
\textbf{Sector} & \textbf{Share of consumption} \\
\midrule
Transport & 40\% \\
Building \& construction & 24\% \\
Packaging & 19\% \\
Engineering & 11\% \\
Consumer durables & 6\% \\
\bottomrule
\end{tabular}
\end{center}

Aluminium serves seven of the fourteen critical industrial ecosystems identified in the EU's
Industrial Strategy, making it a truly cross-cutting material for the green transition.

\subsection{Transport and Mobility}

Aluminium's low density, high strength-to-weight ratio, corrosion resistance, formability, and
recyclability make it indispensable in the automotive and transport sectors. The average aluminium
content per passenger vehicle (Content Per Vehicle, CPV) has risen steadily from 121\,kg in 2006
to \textbf{205\,kg} in 2022, and is projected to reach \textbf{256\,kg} by 2030 (CAGR of 1.9\%
over 2022--2030). Growth is primarily driven by:

\begin{itemize}
    \item \textbf{Castings} (123\,kg/vehicle in 2022): the largest product form, increasingly used
          in battery pack housings, e-drive housings, and mega-castings for body-in-white
          components.
    \item \textbf{Extrusions} (27\,kg/vehicle in 2022): the fastest-growing form, driven by EV
          battery pack housing, body-in-white sills, and electric brake systems (CAGR 5.6\%
          through 2030).
    \item \textbf{Sheet} (45\,kg/vehicle in 2022): growth driven by electrification, closures
          (front and rear doors), and battery cooling plates (CAGR 3.5\% through 2030).
\end{itemize}

\subsection{Construction}

Aluminium's dimensional stability, durability, and recyclability underpin applications that directly
contribute to sustainable buildings: thermally broken window and curtain-wall systems, façade
cladding, roofing, and photovoltaic mounting structures. The material enables natural lighting,
energy savings, and air tightness.

\subsection{Solar Energy}

Aluminium is the most-used metal in solar photovoltaic panels, appearing in frames, support
structures, mounting systems, and electrical wiring. To meet the European Solar Strategy's target
of increasing solar PV capacity by more than 185\% between 2022 and 2030, an estimated
\textbf{10 million tonnes} of aluminium will be required.

\subsection{Packaging}

Aluminium's unique combination of high formability, light weight, strength, and total barrier
properties (against light, gases, and moisture), together with its recyclability, makes it the
preferred packaging material for food and beverages.

\section{Description of Main Processes}

\subsection{Alumina Production (Bayer Process)}

Bauxite ore is crushed, ground, and digested with caustic soda and lime at elevated temperature and
pressure. The resulting alumina-rich solution is clarified, precipitated, filtered, and finally
calcined at high temperature to yield calcined alumina ($\text{Al}_2\text{O}_3$). European plants
consume approximately \textbf{8--10\,GJ of natural gas per tonne of alumina} (equivalent to
$\sim$200\,kg of natural gas), plus 100--200\,kWh of electricity per tonne.

\subsection{Primary Smelting (Hall–Héroult Electrolysis)}

Alumina is dissolved in molten synthetic cryolite ($\text{Na}_3\text{AlF}_6$) to lower the melting
point, and electrolysed in large reduction cells (pots). The overall reaction is:

\[
2\,\text{Al}_2\text{O}_3 + 3\,\text{C} \;\longrightarrow\; 4\,\text{Al} + 3\,\text{CO}_2
\]

Carbon anodes are consumed in the process, generating direct CO$_2$ emissions of approximately
\textbf{1.6--1.8\,tCO$_2$/t Al}. Electric energy consumption is \textbf{14--15\,MWh per tonne}
of aluminium — making the carbon intensity of smelting highly dependent on the electricity
generation mix. A promising breakthrough technology, the \textbf{inert anode} (exemplified by the
Elysis project), replaces carbon anodes with an inert material, producing oxygen instead of CO$_2$
and eliminating direct process emissions at the smelting stage entirely. Indirect emissions from
electricity consumption remain and depend on the grid mix.

\subsection{Semi-Fabrication}

After direct-chill casting into billets or slabs, aluminium is processed into semi-finished products:

\begin{itemize}
    \item \textbf{Rolling}: homogenisation $\rightarrow$ hot rolling (break-down and tandem mills)
          $\rightarrow$ cold rolling $\rightarrow$ annealing, producing sheet and foil.
    \item \textbf{Extrusion}: heated billets are forced through a die to produce profiles of
          virtually any cross-section.
\end{itemize}

\subsection{Recycling}

End-of-life aluminium is shredded, sorted, and remelted in either reverbatory or rotating furnaces
to recover secondary aluminium. Recycling requires only $\sim$8.5\,MJ/kg of primary energy and
$\sim$0.5\,kg\,CO$_2$eq/kg, compared to $\sim$140\,MJ/kg and $\sim$6.7\,kg\,CO$_2$eq/kg for
primary production — an energy saving of approximately \textbf{95\%}.

\section{GHG Emissions Along the Value Chain}

The European aluminium industry (EU + EFTA) emits approximately \textbf{30 Mt\,CO$_2$eq per
year} in total, distributed as follows:

\begin{center}
\begin{tabular}{lcc}
\toprule
\textbf{Segment} & \textbf{Intensity} & \textbf{Total (Mt\,CO$_2$eq/yr)} \\
\midrule
Alumina production    & 0.9\,t\,GHG/t\,Al$_2$O$_3$ & $\sim$6  \\
Smelting              & 4--5\,t\,GHG/t\,Al           & $\sim$17 \\
Semi-fabrication      & 0.5\,t\,GHG/t\,Al            & $\sim$5  \\
Recycling             & 0.5\,t\,GHG/t\,Al            & $\sim$2  \\
\midrule
\textbf{Total}        &                               & $\sim$\textbf{30}  \\
\bottomrule
\end{tabular}
\end{center}

Smelting dominates, with about 50\% of its footprint attributable to indirect electricity-related
emissions. This 30\,Mt figure represents approximately \textbf{5\%} of total European process
industry GHG emissions, and around \textbf{60\%} of the European non-ferrous metals sector.

Crucially, Europe's primary aluminium carbon intensity has improved dramatically: from
$\sim$15\,kg\,CO$_2$eq/kg in 1990 to \textbf{6.6\,kg\,CO$_2$eq/kg} in 2023 — a reduction of
more than 55\% — and is now \textbf{44\% below the global average} of 14.8\,kg\,CO$_2$eq/kg.
Chinese production stands at approximately 20\,kg\,CO$_2$eq/kg. Despite producing 6--7\% of
global primary aluminium, Europe accounts for less than 3\% of the global aluminium industry's
GHG emissions, owing to its relatively clean electricity mix, predominantly hydropower and
nuclear.

\section{Demand Outlook and the Green Transition}

Global aluminium demand is set to increase significantly by 2050, driven by the energy transition.
Under the IAI 2020 reference scenario, total fabricated demand is projected to reach approximately
\textbf{145\,Mt} by 2050, up from around 80\,Mt in 2020. Key demand drivers are electric vehicles
(requiring 59\% more aluminium by 2050 versus 2020 in a 2-degree scenario) and solar PV
(requiring 35\% more). Europe's energy transition alone is expected to require a \textbf{30\%
increase} in aluminium consumption by 2040 compared to 2022 levels.

Within Europe, production forecasts through 2050 anticipate:

\begin{itemize}
    \item Recycling capacity growing by \textbf{+173\%} (remelting and refining), reflecting the
          increasing availability of end-of-life scrap.
    \item Extrusion production growing by \textbf{+18\%} and sheet production by \textbf{+13\%}.
    \item Higher aluminium demand in primary smelting (+22\%) met largely through increased
          imports to compensate for projected European capacity losses.
\end{itemize}

\section{Decarbonisation Pathways}

To achieve net-zero by 2050, three broad decarbonisation routes are identified, drawing on the
European Aluminium \textit{Net-Zero by 2050: Science-Based Decarbonisation Pathways} report
(November 2023):

\begin{enumerate}[label=\arabic*.]
    \item \textbf{Decarbonisation of consumed electricity}: as smelting and increasingly
          semi-fabrication rely on electrical power, a shift to renewable or low-carbon generation
          (wind, solar, hydro, nuclear) is the single most impactful lever. Electricity demand from
          the European aluminium industry is projected to increase by \textbf{+50\,TWh} between
          2020 and 2050, driven mainly by electrification of thermal processes.

    \item \textbf{Decarbonisation of industrial processes}:
    \begin{itemize}
        \item \textit{Inert anode technology} at smelters, eliminating direct CO$_2$ from
              electrolysis.
        \item \textit{Electric boilers and Mechanical Vapour Recompression (MVR)} for steam
              generation in alumina refining (Bayer process).
        \item \textit{Electric furnaces} for calcination, melting, and annealing in both
              semi-fabrication and recycling.
        \item \textit{Green hydrogen} as a fuel substitute, particularly for high-temperature
              processes such as calcination and anode baking; full fuel substitution across the
              European industry could require up to \textbf{1\,Mt\,H$_2$/year}.
        \item \textit{Carbon capture, utilisation, and storage (CCU/CCS)}, especially for
              residual process emissions in alumina calcination.
    \end{itemize}

    \item \textbf{Maximising circularity}: increasing post-consumer scrap collection rates,
          implementing deposit-return schemes (DRS), improving dismantling and sorting
          technologies, and developing smarter alloying strategies to enable closed-loop recycling
          with minimal quality loss.
\end{enumerate}

The decarbonisation pathway modelled by Ramboll shows that total industry emissions could fall from
approximately 20\,Mt\,CO$_2$eq in 2021 to near zero by 2050, with the steepest reductions
occurring between 2025 and 2040, driven first by grid decarbonisation, then by process
electrification, and finally by the shift to inert anode technology.

\section{Case Study: Aluminium Dunkerque}

Aluminium Dunkerque is one of Europe's largest primary aluminium smelters, producing
\textbf{300{,}000 tonnes of primary aluminium per year}. Located in northern France adjacent to
the Gravelines nuclear power station — itself one of the largest in Western Europe with six
910\,MW pressurised water reactors and an annual output of 31--37\,TWh — the smelter draws
primarily on low-carbon nuclear electricity.

Key energy figures for Aluminium Dunkerque:

\begin{itemize}
    \item Specific electricity consumption: \textbf{15\,MWh/tonne of aluminium}.
    \item Total annual electricity consumption: $15 \times 300{,}000 = $\textbf{4.5\,TWh/year}.
    \item Typical continuous power draw: \textbf{514\,MW}.
\end{itemize}

For context, 4.5\,TWh/year is equivalent to three times the annual consumption of all SNCF
(French national railways), and roughly 65\% of the city of Marseille's annual electricity use. In
terms of renewable generation equivalents, it would require approximately 700 wind turbines
(assuming 3\,MW capacity each producing 7{,}000\,MWh/year) or 18 million m$^2$ of photovoltaic
panels (at 250\,kWh/m$^2$/year).

The plant operates two 850-metre-long electrolysis halls housing \textbf{264 pots}, and a foundry
with four casting machines and seven furnaces producing 80\% slabs (for the rolling industry,
packaging, transport, and construction) and 20\% ingots (for the automotive sector).

\section{Conclusions}

Aluminium is an irreplaceable enabler of the green energy transition. Its unique physical properties
— light weight, strength, corrosion resistance, formability, conductivity, and infinite
recyclability — allow key sectors including transport, construction, renewable energy, and packaging
to decarbonise in ways that would otherwise be impossible.

However, primary aluminium production remains one of the most electricity- and GHG-intensive
industrial processes. Meeting the 2050 climate targets requires drastic reductions in GHG intensity
across the entire value chain. The three main routes — electricity decarbonisation, process
electrification and fuel switching, and maximised circularity — are technically available and
increasingly cost-competitive. Policy frameworks that ensure access to affordable low-carbon
electricity and protect European producers from carbon leakage (via the Carbon Border Adjustment
Mechanism) will be essential to enable the industry's transformation while maintaining its
strategic role in the European economy.